\documentclass[12pt]{article}

\usepackage{fouriernc}
\usepackage{amsmath}
\usepackage{amssymb}
\usepackage{pifont}
\newcommand{\cmark}{\text{\ding{51}}}
\newcommand{\xmark}{\text{\ding{55}}}

\usepackage{extarrows}
\usepackage[shortlabels]{enumitem}

\usepackage{tikz}

\newcommand{\contradiction}{%
\begin{tikzpicture}[rotate=45,x=0.5ex,y=0.5ex]
\draw[line width=.2ex] (0,2) -- (3,2) (0,1) -- (3,1) (1,3) -- (1,0) (2,3) -- (2,0);
\end{tikzpicture}
}

\begin{document}

\section*{Problem-225}
\subsection*{Problem Statement}
If in the equation
\[
	ax^2 + bx + c = 0
\]
the coefficient $a$ is equal to zero then the equation takes the form
\[
	bx + c = 0
\]
and has the unique solution
\[
	x = -\frac{c}{b}
\]

Strictly speaking, the last sentence is wrong; when $b=0$ the quotient $\frac{c}{b}$ is undefined. How are we to correct this error?

\subsection*{Solution}
We should require that in the equation $ax^2+bx+c = 0$, both $a$ and $b$ cannot be zero.

\section*{Problem-227}
\subsection*{Problem Statement}
Prove that $\sqrt{3}$ is irrational.

\subsection*{Solution}
The statement `$\sqrt{3}$ is irrational' is equivalent to the statement: $\sqrt{3} \neq \frac{a}{b}$ for any integer $a$ and $b$. We shall make use of \textbf{proof by contradiction}, therefore, we shall assume that $\sqrt{3} = \frac{a}{b}$ and that should lead to a contradiction, completing the proof.

We shall consider four different cases for the pair $(a, b)$ and we shall derive a contradiction for each.
\begin{enumerate}
\item $a$ and $b$ both are odd integer. Say $a = 2m+1$ and $b = 2n+1$, where $m$ and $n$ are integers. According to our assumption
\begin{equation*}
	\begin{aligned}
		\sqrt{3} &= \frac{2m+1}{2n+1}\\
		3        &= \frac{ (2m+1)^2 }{ (2n+1)^2 }\\
		3\ (2n+1)^2 &= (2m+1)^2\\
		3\ (4n^2 + 4n + 1) &= 4m^2 + 4m + 1\\
		12n^2 + 12n + 3 &= 4m^2 + 4m + 1\\
		12n^2 + 12n + 2 &= 4m^2 + 4m\\
		6n^2 + 6n + 1 &= 2m^2 + 2m
	\end{aligned}
\end{equation*}
On the left side we have an odd integer and on the right side we have an even integer---a contradiction. \contradiction

\item $a$ is an odd integer and $b$ is an even integer. Say $a = 2m+1$ and $b = 2n$, where $m$ and $n$ are integers. According to our assumption
\begin{equation*}
	\begin{aligned}
		\sqrt{3} &= \frac{ 2m+1 }{ 2n }\\
		3 &= \frac{ (2m+1)^2 }{ 4n^2 }\\
		12n^2 &= (2m+1)^2
	\end{aligned}
\end{equation*}
On the left side we have an even integer and on the right side we have an odd integer---a contradiction. \contradiction

\item $a$ is an even integer and $b$ is an odd integer. Say $a = 2m$ and $b = 2n+1$, where $m$ and $n$ are integers. According to our assumption
\begin{equation*}
	\begin{aligned}
		\sqrt{3} &= \frac{ 2m }{ 2n+1 }\\
		3 &= \frac{ 4m^2 }{ (2n+1)^2 }\\
		3\ (2n+1)^2 &= 4m^2
	\end{aligned}
\end{equation*}
On the left side we have an odd integer and on the right side we have an even integer---a contradiction. \contradiction

\item $a$ and $b$ both are even integer. In this case, $a$ and $b$ must have 2 as a common factor. We can divide both $a$ and $b$ by 2. We can keep dividing by 2 as long as both remain even integer. At the end, we are in one of the previous three cases.
\end{enumerate}

\section*{Problem-229}
\subsection*{Problem Statement}
Prove that for $a \geq 0,\ b > 0$
\[
	\sqrt{ \frac{a}{b} } = \frac{ \sqrt{a} }{ \sqrt{b} }
\]

\subsection*{Solution}
Basically we are being asked to show that $\frac{ \sqrt{a} }{ \sqrt{b} }$ is the square root of the non-negative integer $\frac{a}{b}$. So, we need to prove the below two:
\begin{itemize}
	\item $\frac{ \sqrt{a} }{ \sqrt{b} }$ is non-negative.\\
	
	Since $a$ is non-negative, its square root must also be non-negative. Since $b$ is positive, its square root must also be positive. $\frac{ \sqrt{a} }{ \sqrt{b} }$ is thus the ratio of a non-negative number and a positive number; so, it must be non-negative.
	\item Squaring $\frac{ \sqrt{a} }{ \sqrt{b} }$ gives us $\frac{a}{b}$.

	\begin{equation*}
			\left( \frac{ \sqrt{a} }{ \sqrt{b} } \right)^2 = \frac{ (\sqrt{a})^2 }{ (\sqrt{b})^2 } = \frac{a}{b}.
	\end{equation*}		
\end{itemize}

\section*{Problem-230}
\subsection*{Problem Statement}
Is the equality $\sqrt{ a^2 } = a$ true for all $a$?

\subsection*{Solution}
It might seem counter-intuitive that the answer is no. But if we observe closely, the question is actually asking about the equality $+\sqrt{ a^2 } = a$, where the `+' is implied. When we make the `+' explicit, it is more noticeable that for $a < 0$, the equality cannot hold, because the left side of the equality is always positive. The  equality that works for all $a$ is $\sqrt{ a^2 } = |a|$, where
\[
	|a| = 
		\begin{cases}
		   \ \ a &\mbox{, if } a \geq 0\\
		   -a &\mbox{, if } a < 0
		\end{cases}
\]

It is unfortunate that we do not always make the `+' explicit.

\section*{Problem-232}
\subsection*{Problem Statement}
Which is bigger: $\sqrt{1001} - \sqrt{1000}$ or $\frac{1}{10}$?
\subsection*{Solution}
It would be easier to compare if we could convert $\sqrt{1001} - \sqrt{1000}$ into a fraction.
\begin{equation*}
	\begin{aligned}
		& \sqrt{1001} - \sqrt{1000}\\
		&= \frac{ \left(\sqrt{1001} - \sqrt{1000}\right)\ \left(\sqrt{1001} + \sqrt{1000}\right) }{ \sqrt{1001} + \sqrt{1000} }\\
		&= \frac{ \left( \sqrt{1001} \right)^2 - \left( \sqrt{1000} \right)^2 }{ \sqrt{1001} + \sqrt{1000} }\\
		&= \frac{1001 - 1000}{ \sqrt{1001} + \sqrt{1000} }\\
		&= \frac{1}{ \sqrt{1001} + \sqrt{1000} }\\
		&< \frac{1}{ \sqrt{900} + \sqrt{900} }
	\end{aligned}
\end{equation*}
Since $\sqrt{900} + \sqrt{900} = 60$, $\sqrt{1001} - \sqrt{1000} < \frac{1}{60} < \frac{1}{10}$. So, $\frac{1}{10}$ is the bigger between the two.

\section*{Problem-239}
\subsection*{Problem Statement}
\subsubsection{Solutions to quadratic equation}
A quadratic equation $x^2+px+q = 0$ can be solved by ``completing the square''.
\begin{equation*}
	\begin{aligned}
		x^2 + 2 \cdot x \cdot \frac{p}{2} + \left( \frac{p}{2} \right)^2 + \left( q - \frac{p^2}{4} \right) &= 0\\
		\left( x+\frac{p}{2} \right)^2 &= \frac{p^2}{4} - q\\
		x &= -\frac{p}{2} + \frac{ \pm \sqrt{p^2 - 4q} }{2}
	\end{aligned}
\end{equation*}
There are three possible cases:
\begin{enumerate}
	\item $p^2-4q > 0$. In this case two distinct solutions $x_1,\ x_2$ exist: 
	
	$x_1 = -\frac{p}{2} + \frac{ \sqrt{ p^2-4q } }{2}$ and $x_2 = -\frac{p}{2} - \frac{ \sqrt{ p^2-4q } }{2}$
	\item $p^2-4q = 0$. In this case the two solutions coincide; therefore, there is one distinct solution.
	
	$x_1 = -\frac{p}{2}$ and $x_2 = -\frac{p}{2}$
	\item $p^2-4q < 0$. In this case there is no real-valued solution.
\end{enumerate}
Since the sign of $p^2-4q$ determines how many distinct solutions there are, for convenience, we shall call this expression $D$, therefore, $D = p^2-4q$.

\subsubsection{Vieta's theorem}
If a quadratic equation $x^2 + px + q = 0$, has two distinct solutions $\alpha$ and $\beta$ then
\[
	\alpha + \beta = -p
\]
\[
	\alpha \cdot \beta = q
\]

\subsubsection{Proof-1 of Vieta's theorem}
If $x^2 + px + q = 0$ has two distinct solutions, we are in the first case of $0.0.1$; thus we can let $\alpha$ and $\beta$ be as follows:
\setcounter{equation}{0}
\begin{equation}
	\alpha = -\frac{p}{2} + \frac{ \sqrt{ D } }{2}
\end{equation}
\begin{equation}
	\beta = -\frac{p}{2} - \frac{ \sqrt{ D } }{2}
\end{equation}
$(1)+(2)$ gives:
\[
	\alpha + \beta = -p
\]
$(1) \times (2)$ gives:
\begin{equation*}
	\begin{aligned}
		\alpha \cdot \beta &= \frac{p^2}{4} - \frac{D}{4}\\
						   &= \frac{ p^2 - \left( p^2-4q \right) }{4}\\
						   &= q
	\end{aligned}
\end{equation*}

\subsubsection{Proof-2 of Vieta's theorem}
We can rewrite $x^2 + px + q = 0$ as $P(x) = 0$ where $P(x) = x^2 + px + q$ is a polynomial with degree two. If $\alpha$ and $\beta$ are two distinct solutions of $P(x) = 0$ then
\[
	P(x) = (x-\alpha)\ (x-\beta)\ R(x)
\]
Here $R(x)$ must be constant, otherwise right side would have degree more than two. In fact $R(x)$ must be exactly 1, otherwise the coefficient of $x^2$ would not match with the one in $P(x)$. So, we have
\begin{equation*}
	\begin{aligned}
		P(x) &= (x - \alpha)\ (x - \beta)\\
			 &= x^2 -(\alpha + \beta)\ x + \alpha \cdot \beta
	\end{aligned}
\end{equation*}
Comparing $P(x) = x^2 + (\ -(\alpha+\beta)\ )\ x + \alpha \cdot \beta$ with $P(x) = x^2 + px + q$ gives us Vieta's theorem.

\subsubsection{Questions to answer}
Can you generalize Vieta's theorem to the case of a quadratic equation having only one root? Are both proofs still valid for this case?

\subsection{Solution}
\subsubsection{Generalization}
Generalized version of Vieta's theorem is:\\

If a quadratic equation $x^2 + px + q = 0$ has two roots $\alpha$ and $\beta$ (the roots can be equal) then
\[
	\alpha + \beta = -p
\]
\[
	\alpha \cdot \beta = q
\]

\subsubsection{Proof-1}
Vieta's generalized theorem already works for the first case of $0.0.1$. We shall show that it works for the second case as well. In the second case of $0.0.1$, $\alpha = \beta = -\frac{p}{2}$.
\[
	\alpha + \beta = 2\alpha = 2 \left( -\frac{p}{2} \right) = -p
\]
In the second case of $0.0.1$ we also have $p^2-4q = 0$; thus $q = \frac{p^2}{4}$.
\[
	\alpha \cdot \beta = \alpha^2 = \frac{p^2}{4} = q
\]

\subsubsection{Proof-2}
Since $x^2+px+q = 0$ has two duplicate solutions, we can consider $P(x) = x^2 + px + q$ having two repeated roots, each equal to $\alpha$. Thus $P(x) = (x-\alpha)\ (x-\alpha)$. In other words, $P(x) = x^2 - 2\alpha + \alpha^2$. A comparison gives us Vieta's theorem.


\section*{Problem-240}
\subsection*{Problem Statement}
(Vieta's theorem for a cubic equation) Assume that a cubic equation $x^3 + px^2 + qx + r = 0$ has three different roots $\alpha,\ \beta,\ \gamma$. Prove that
\begin{equation*}
	\begin{aligned}
		\alpha + \beta + \gamma &= -p\\
		\alpha \beta + \alpha \gamma + \beta \gamma &= q\\
		\alpha \beta \gamma &= -r
	\end{aligned}
\end{equation*}

\subsection{Solution}
The polynomial $P(x) = x^3 + px^2 + qx + r$ can be written as
\[
	P(x) = (x - \alpha)\ (x - \beta)\ (x - \gamma)\ R(x)
\]
To match the coefficient of $x^3$, $R(x)$ must be 1. Therefore,
\begin{equation*}
		\begin{aligned}
			P(x) &= (x - \alpha)\ (x - \beta)\ (x - \gamma)\\
				 &= x^3 - (\alpha + \beta + \gamma)\ x^2 + ( \alpha \cdot \beta + \alpha \cdot \gamma + \beta \cdot \gamma )\ x - \alpha\cdot \beta\cdot \gamma
		\end{aligned}
\end{equation*}
On comparing against $P(x) = x^3 + px^2 + qx + r$ we have
\begin{equation*}
	\begin{aligned}
		\alpha + \beta + \gamma &= -p\\
		\alpha \beta + \alpha \gamma + \beta \gamma &= q\\
		\alpha \beta \gamma &= -r
	\end{aligned}
\end{equation*}

\section*{Problem-259}
\subsection*{Problem Statement}
Starting with the graph for $y = x^2$ we can get the graph for $y = a(x+m)^2 + n$ as follows:
\begin{enumerate}[label=(\alph*)]
	\item Stretching it vertically $a$ times gives us the graph for $y = ax^2$. If $a > 1$, the graph gets steeper. If $0 < a < 1$, the graph gets flatter. If $a < 0$, the graph becomes upside down.
	\item Moving the resulting graph $m$ units to the left gives us the graph for $y = a(x+m)^2$. If $m > 0$, the graph moves to the left. If $m < 0$, the graph moves to the right.
	\item Moving the resulting graph $n$ units up gives us the graph for $y = a(x+m)^2 + n$. If $n > 0$, the graph moves upwards. If $n < 0$, the grpah moves downwards.
\end{enumerate}

There are six possible orderings of operations (a), (b), and (c). Do we get six different graphs or do some of the graphs coincide?

\subsection*{Solution}
We do not get six different graphs. The six graphs fall into two distinct groups: one for the function $y = a(x+m)^2+n$ and the other for the function $y = a(x+m)^2+an$. Since (b) is the only one that affects $x-$coordinates, (b) does not interfere with (a) and (c). On the other hand, (a) and (c) both affect $y-$coordinates and thus they interfere with each other. Which of these two groups a graph falls into, thus, depends on the order of (a) and (c). If (c) follows (a), we get the graph for the function $y = a(x+m)^2+n$. If (a) follows (c), we get the graph for the function $y = a(x+m)^2+an$.

\begin{enumerate}[1.]
	\item (c) follows (a):
		\begin{enumerate}[i.]
			\item $y=x^2 \xLongrightarrow{\text{(a)}} y = ax^2 \xLongrightarrow{\text{(b)}} y = a(x+m)^2 \xLongrightarrow{\text{(c)}} y = a(x+m)^2+n$
			\item $y=x^2 \xLongrightarrow{\text{(a)}} y=ax^2 \xLongrightarrow{\text{(c)}} y=ax^2 + n \xLongrightarrow{\text{(b)}} y=a(x+m)^2+n$
			\item $y=x^2 \xLongrightarrow{\text{(b)}} y=(x+m)^2 \xLongrightarrow{\text{(a)}} y=a(x+m)^2 \xLongrightarrow{\text{(c)}} y=a(x+m)^2+n$	
		\end{enumerate}
	\item (a) follows (c):
		\begin{enumerate}[i.]
			\item $y=x^2 \xLongrightarrow{\text{(b)}} y=(x+m)^2 \xLongrightarrow{\text{(c)}} y=(x+m)^2+n \xLongrightarrow{\text{(a)}} y=a(x+m)^2+an$
			\item $y=x^2 \xLongrightarrow{\text{(c)}} y=x^2+n \xLongrightarrow{\text{(a)}} y=ax^2+an \xLongrightarrow{\text{(b)}} y=a(x+m)^2+an$
			\item $y=x^2 \xLongrightarrow{\text{(c)}} y=x^2+n \xLongrightarrow{\text{(b)}} y=(x+m)^2+n \xLongrightarrow{\text{(a)}} y=a(x+m)^2+an$
		\end{enumerate}
\end{enumerate}

\section*{Problem-260}
\subsection*{Problem Statement}
How can you determine the signs of $a,\ b,\ c$ by looking at the graph of $y = ax^2+bx+c$?

\subsection{Solution}
If $a = 0$ but $b \neq 0$, we have $y = bx+c$, thus our graph is no longer a parabola, but a straight line. The slope and y-intercept of the straight line gives the signs of $b$ and $c$. On the other hand, if $b = 0$ but $a \neq 0$, we have $y = ax^2 + c$. This is similar to the form $y = a(x+m)^2+n$, so we would know the signs of $a$ and $c$ easily. If both $a$ and $b$ are zero, we have $y = c$. So, we have a graph that is a straight line and parallel to the $x-$axis. So, if the line lies above the $x-$axis, $c > 0$, else $c < 0$. 
We shall now consider the more interesting case where $a \neq 0$ and $b \neq 0$. Since we know a lot about the graph of $y = a(x+m)^2+n$, let's transform $y = ax^2 + bx + c$ into that form.

\begin{equation*}
	\begin{aligned}
		y &= ax^2 + bx + c\\
		  &= a\left( x^2 + x\cdot \frac{b}{a} \right) + c\\
		  &= a\left( x^2 + 2 \cdot x \cdot \frac{b}{2a} + \frac{b^2}{4a^2} \right) + c - \frac{b^2}{4a^2}\\
		  &= a \left( x+\frac{b}{2a} \right)^2 + \left( c - \frac{b^2}{4a^2} \right)
	\end{aligned}
\end{equation*}

Comparing with $y = a(x+m)^2 + n$, we can consider $\frac{b}{2a} = m$.

Below is how we can find out the signs of $a,\ b,\ c$ :
\begin{itemize}
	\item Compared to the graph of $y = x^2$, if our graph is upside down, then $a < 0$; otherwise $a > 0$.
	\item $\frac{b}{2a}$ determines how much our graph is shifted horizontally compared to the graph of $y = x^2$. So, if we know the sign of $a$ and how the graph is shifted horizontally, then we can deduce the sign of $b$. 
	
\begin{table}[ht]
\centering % used for centering table
\begin{tabular}{| l | c | c |} % centered columns (4 columns)
\hline %inserts double horizontal lines
 & $a > 0$ & $a < 0$ \\ % inserts table
%heading
\hline % inserts single horizontal line
Shifted left  & $b > 0$ & $b < 0$ \\
Shifted right & $b < 0$ & $b > 0$ \\
Unshifted     & $b = 0$ & $b = 0$ \\
\hline %inserts single line
\end{tabular}
\end{table}

	\item When $x = 0$, $y = ax^2 + bx + c$ becomes $y = c$. So, looking at where our graph intersects with the line $x = 0$ aka the $y-$axis, we would know the sign of $c$.
\end{itemize}

\section*{Problem-262}
\subsection*{Problem Statement}
The sum of two numbers is equal to 1. What is the maximal possible value of their product?

\subsection*{Solution}
We see that 347 + (-346) = 1, but their product is negative. Whereas 0 + 1 = 1, and their product is 0 and this is better. Can the product be positive? What about $\frac{1}{10} + \frac{9}{10}$? In this case the product is $\frac{9}{100}$ and is the best among the three examples we have just considered.

Say one number is $x$ and the other is $1-x$. So, the product is the polynomial $x\ (1-x) = -x^2+x = -\left( x - \frac{1}{2} \right)^2 + \frac{1}{4}$. So, the graph of the product polynomial is a parabola that is upside down, $\frac{1}{2}$ units shifted right; therefore the vertex or crest is at $x = \frac{1}{2}$ and that is where maximal value occurs. So, the maximal value is $\frac{1}{2}\ \left( 1 - \frac{1}{2} \right) = \frac{1}{4}$.

\section*{Problem-263}
\subsection*{Problem Statement}
Prove that a square has the maximum area of all rectangles having the same perimeter.

\subsection*{Solution}
Say the perimeter is $p$. If one side of the rectangle is $x$, then the other side must be $\frac{p}{2} - x$. Now the area is a polynomial:  $x\ \left( \frac{p}{2}-x \right) = -x^2 + x\cdot \frac{p}{2} = -\left( x - \frac{p}{4} \right)^2 + \frac{p^2}{16}$. Thus the area has the maximum value when $x = \frac{p}{4}$. But in that case the other side is also $\frac{p}{4}$, making it a square.

\section*{Problem-264}
\subsection*{Problem Statement}
Prove that a square has the minimum perimeter of all rectangles having the same area.

\subsection*{Solution}
So we need to prove that among all rectangles with an area say $a$, the square has the minimum perimeter. Say $x$ is one side of a rectangle with area $a$. The other side then is $\frac{a}{x}$. The perimeter is a polynomial $P(x) = 2\ \left( x + \frac{a}{x} \right)$. Consider the below equation, where $c > 0$.
\setcounter{equation}{0}
\begin{equation}
	2\ \left( x + \frac{a}{x} \right) = c
\end{equation}
We shall first find out for which $c > 0$, the equation $(1)$ has solutions. Then from those values of $c > 0$, we shall find out the minimum, because that gives us the minimum periemeter for a given area $a$. We shall then see, for that minimum $c > 0$, what value $x$ takes. Once we know how $x$ and the other side $\frac{a}{x}$ look, we would know what type of rectangle has this minimum perimeter. Note, since this is a rectangle we are talking about, we are interested in $x > 0$ solutions.

\begin{equation*}
	\begin{aligned}
		2\ \left( x + \frac{a}{x} \right) &= c\\
		x^2 - x\cdot \frac{c}{2} + a &= 0\\
		\left( x - \frac{c}{4} \right)^2 &= \frac{c^2}{16} - a\\
		x &= \frac{c}{4} + \frac{ \pm \sqrt{ c^2 - 16a } }{16}
	\end{aligned}
\end{equation*}

So, $(1)$ has solutions when $c^2-16a \geq 0$. In other words, when $c \geq 4\sqrt{a}$ or $c \leq -4\sqrt{a}$, the equation $(1)$ has solutions. Hence, the minimum positive value of $c$ for which $(1)$ has a solution is $4\sqrt{a}$ and at this value of $c$, we have $x = \frac{c}{4} = \frac{4\sqrt{a}}{4} = \sqrt{a}$. The other side is $\frac{a}{x} = \frac{a}{\sqrt{a}} = \sqrt{a}$. Therefore, the rectangle with area $a$ that has the minimum perimeter is a square.\\

In another way, say a square with area $a$ has perimeter $p$. Now, from \textbf{Problem-263}, we know that $a$ is the maximum possible area with the perimeter $p$. For all other rectangles, if they want to achieve the area $a$, their perimeter must be bigger than $p$. So, with area $a$, $p$ is the minimum possible perimeter for a rectangle and that happens when the rectangle is a square.

\section*{Problem-265}
\subsection*{Problem Statement}
Find the minimal value of the expression $x + \frac{2}{x}$ for positive $x$.

\subsection*{Solution}
We may consider $x$ and $\frac{2}{x}$ as the two sides of a rectangle with area 2. With that interpretation, $2\ \left( x + \frac{2}{x} \right)$ is the perimeter of the rectangle which, according to \textbf{Problem-264}, is minimal when the rectangle is a square. Therefore, $2\ \left( x + \frac{2}{x} \right)$ is minimal when $x = \frac{2}{x}$. This is also the condition for our expression $x + \frac{2}{x}$, which is half the perimeter, to have its minimum value for $x > 0$. So, when $x = \sqrt{2}$, the perimeter and also our expression $x + \frac{2}{x}$ assume their minimal values. So minimal value of $x + \frac{2}{x}$ is $\sqrt{2} + \frac{2}{\sqrt{2}} = 2\sqrt{2}$.

\section*{Problem-267}
\subsection*{Problem Statement}
Construct a biquadratic equation
\begin{enumerate}[(a)]
	\item having no solution;
	\item having exactly one solution;
	\item having exactly two solutions;
	\item having exactly three solutions;
	\item having exactly four solutions;
	\item having exactly five solutions;
\end{enumerate}

\subsection*{Solution}
A biquadratic equation has the form $ax^4+bx^2+c = 0$. We are seeking solution in real numbers.
\begin{enumerate}[(a)]
	\item $x^4+2x^2+1 = 0$ does not have a solution. Because if there were a solution then $y = x^2$ would have been a solution to the equation $y^2+2y+1 = 0$. Then $(y+1)^2 = 0$, or $y = -1$. That means $x^2 = -1$.
	\item $x^2(x^2+1) = 0$ or $x^4+x^2 = 0$ has exactly one solution $x = 0$.
	\item $\left(x^2-1\right)\ \left(x^2+1\right) = 0$ or $x^4-1 = 0$ has exactly two solutions, namely $x_{1,2} = \pm 1$.
	\item We modify the solution for (b) to get a biquadratic equation having exactly three roots. $x^2\ \left(x^2-1\right) = 0$ or $x^4-x^2 = 0$ has exactly three roots: $x_{1,2,3} = 0, \pm 1$.
	\item $\left(x^2-1\right)\ \left(x^2-2\right) = 0$ or $x^4-3x^2+2=0$ has exactly four solutions, namely $x_{1,2,3,4} = \pm 1, \pm \sqrt{2}$.
	\item A biquadratic equation cannot have exactly five solutions. Because if there were such solutions, we could write the equation like below:
	\[
		(x-a)\ (x-b)\ (x-c)\ (x-d)\ (x-e) = 0
	\]
And the highest degree of the polynomial on the left would then be more than 4 and it would not be a biquadratic equation anymore.
\end{enumerate}

\section*{Problem-268}
\subsection*{Problem Statement}
What is the possible number of solutions of the equation
\[
	ax^6 + bx^3 + c = 0\ ?
\]

\subsection*{Solution}
If $x$ is a solution to our equation, then $y = x^3$ is a solution to the quadratic equation $ay^2 + by + c = 0$. As long as $a \neq 0$, the quadratic equation can have $0,\ 1,$ or $2$ solutions. There will also be 0, 1, or 2 corresponding solutions to our equation $ax^6 + bx^3 + c = 0$. If $a = 0$ and $b \neq 0$, we have exactly one solution. If $a = 0, b = 0, c = 0$, we have infinitely many solutions.

\section*{Problem-269}
\subsection*{Problem Statement}
The same question for the equation
\[
	ax^8 + bx^4 + c = 0
\]

\subsection*{Solution}
There are 0, 1, 2, 3, 4, or infinitely many solutions. Below table shows the different number of solutions depending on which of $a,\ b,\ c$ are non-zero.

\begin{table}[ht]
\centering % used for centering table
\begin{tabular}{| r | r | r | r |} % centered columns (4 columns)
\hline %inserts double horizontal lines
 $a = 0$ & $b = 0$ & $c = 0$ & Maximum Number of Solutions \\ % inserts table
%heading
\hline % inserts single horizontal line
\hline
$\cmark$ & $\cmark$ & $\cmark$ & $\infty$ \\
\hline
$\cmark$ & $\cmark$ & $\xmark$ & 0\\
\hline
$\cmark$ & $\xmark$ & $\cmark$ & 1, ($x = 0$)\\
\hline
$\cmark$ & $\xmark$ & $\xmark$ & 2, $\left( x_{1,2} = {\pm \beta}^{\frac{1}{4}} \text{ with } \beta = -\frac{c}{b} > 0 \right)$\\
\hline
$\xmark$ & $\cmark$ & $\cmark$ & 1, ($x=0$)\\
\hline
$\xmark$ & $\cmark$ & $\xmark$ & 2, $\left( x_{1,2} = \pm \gamma^{\frac{1}{8}} \mbox{ with } \gamma = -\frac{c}{a} > 0 \right)$\\
\hline
$\xmark$ & $\xmark$ & $\cmark$ & 3, $\left( x_{1,2,3}=0, \pm \delta^{\frac{1}{4}} \mbox{ with } \delta = -\frac{b}{a} > 0 \right)$\\
\hline
$\xmark$ & $\xmark$ & $\xmark$ & \shortstack{4, two come from each of the two positive solutions of \\ $ay^2 + ay+c = 0$ with $y = x^4$}\\
\hline %inserts single line
\end{tabular}
\end{table}

\section*{Problem-270}
\subsection*{Problem Statement}
Solve the equation
\[
	2x^4 + 7x^3 + 4x^2 + 7x + 2 = 0
\]

\subsection{Solution}
This is a symmetric equation with a center at $4x^2$, which means the powers of $x$ with same coefficient left and right of the center are equally distant from the center-power. For example, $x^4$ on the left of $x^2$ and $x^0$ on the right of the $x^2$ both have the same coefficient 2. So the idea would be to distribute a bit of the power from the left to the right so that we have a matching pair of powers for $x$, albeit with opposing signs. To do that, we shall divide the equation by $x^2$. We can do that because when $x=0$, the equation reduces to $2 = 0$, thus $x = 0$ isn't a solution of the equation.
\begin{equation*}
	\begin{aligned}
		2x^2 + 7x + 4 + \frac{7}{x} + \frac{2}{x^2} &= 0\\
		2\left( x^2 + \frac{1}{x^2} \right) + 7\ \left( x + \frac{1}{x} \right) + 4 &= 0\\
	\end{aligned}
\end{equation*}
We note that $\left( x + \frac{1}{x} \right)^2 = x^2 + \frac{1}{x^2} + 2$, therefore, $x^2 + \frac{1}{x^2} = \left(x + \frac{1}{x}\right)^2-2$. Using this observation in the above equation, we have:
\begin{equation*}
	\begin{aligned}
		2\ \left(x + \frac{1}{x}\right)^2 + 7\ \left(x + \frac{1}{x}\right) &= 0
	\end{aligned}
\end{equation*}
Let $x + \frac{1}{x}$ be called $u$, for convenience.
\begin{equation*}
	\begin{aligned}
		2u^2 + 7u &= 0\\
		u \cdot (2u+7) &= 0\\
		2u \cdot \left(u+\frac{7}{2}\right) &= 0
	\end{aligned}
\end{equation*}
So $u = 0$ or $u = -\frac{7}{2}$. If $u = 0$, we end up having $x^2+1 = 0$, so that one does not give a solution. On the other hand, $u = -\frac{7}{2}$ gives us $2x^2 + 7x + 2 = 0$ and that gives us two solutions: $x_{1,2} = \frac{-7 \pm \sqrt{33}}{4}$.

\section*{Problem-272}
\subsection*{Problem Statement}
Compute $\sqrt[7]{0.999}$ to three decimal places.

\subsection*{Solution}
Let's see where $\sqrt[7]{0.999}$ lies with respect to $0.9994$. Here $0.9994$ is interesting, because we know $\sqrt[7]{0.999} < 1$; so, if the required 7-th root lies above $0.9994$ we can say that $0.9995 \leq \sqrt[7]{0.999} < 1$; therefore, rounded to three decimal places, the 7-th root would be $1.000$. On the other hand, if $0.9994 > \sqrt[7]{0.999}$ we need to do more work. To find out if $0.9994 < \sqrt[7]{0.999}$ we shall see if $0.9994^7 < 0.999$ is true. Let's consider $0.9994^2$ as a starter.
\begin{equation*}
	\begin{aligned}
		0.9994^2 &= (1-0.0006)^2\\
		         &= 1^2 - 2 \times 1 \times 6 \times 10^{-4} + \left( 6 \times 10^{-4} \right)^2\\
		         &= 1 - 12 \times 10^{-4} + 36 \times 10^{-8}\\
		         &= 0.99880036\\
		         &< 0.999
	\end{aligned}
\end{equation*}
We see that even the second power of 0.9994 is less than 0.999. The seventh power of 0.9994 is even smaller. Thus $0.9994 < \sqrt[7]{0.999}$ and to three decimal places $\sqrt[7]{0.999} = 1.000$.

\end{document}




















